\section{Higher extensions of the pro-generalized Jacobian}\label{sec:H}
\subsection{The exact category and the positive result}
Throughout this section $k$ is algebraically closed of characteristic $p>0$,
and $\mathcal P$ is Serre's abelian category of commutative
pro-quasi-algebraic groups. Algebraic groups are viewed through their
quasi-algebraic perfection. We do not identify this category with the
category of all sheaves on a perfect \emph{big} \'etale site.
For a smooth curve with function field $K$, use its smooth projective
completion $C$. Write $J_{\mathfrak n}$ for the connected degree-zero
generalized Jacobian of a nonzero effective modulus $\mathfrak n$.

\begin{theorem}[H1: higher Ext vanishing; Proved]\label{h:vanishing}
The object $J_\infty=\varprojlim_{\mathfrak n}J_{\mathfrak n}$ has
projective dimension at most one in $\mathcal P$. Thus
\[
 \Ext^i_{\mathcal P}(J_\infty,G)=0\qquad(i\geq2)
\]
for every $G\in\mathcal P$, including all the algebraic targets in Part H
and characteristics two and three.
\end{theorem}
\begin{proof}
For $\mathfrak n=\sum_{x\in S}n_x[x]$, put
\[
 R_{\mathfrak n}=\prod_{x\in S}
  (\OO_{C,x}/\mathfrak m_x^{n_x})^\times,
 \qquad L_{\mathfrak n}=R_{\mathfrak n}/\GG_m,
\]
with diagonal scalar units. The generalized-Jacobian sequence is
\[
 0\longrightarrow L_{\mathfrak n}\longrightarrow J_{\mathfrak n}
 \longrightarrow J=\Pic^0(C)\longrightarrow0.
\]
Its algebraic-group assertion is \cite[Chapter I, \S1, Theorem 3,
p.~12]{SerreGACC}. Equivalently, rigidifications of a degree-zero line
bundle along $\mathfrak n$ give this sequence; rigidifications of the
trivial bundle are $R_{\mathfrak n}$ modulo its scalar automorphisms.
Forgetting rigidifications gives compatible transition morphisms.
Inverse limits are exact in $\mathcal P$
\cite[\S2.5, Proposition 10, p.~19]{Serre1960}, so
\[
 0\longrightarrow L\longrightarrow J_\infty\longrightarrow J
 \longrightarrow0,\qquad L=\varprojlim L_{\mathfrak n}.
\]
Choose $x_0\in C(k)$ and restrict to the cofinal moduli containing it.
Residue at $x_0$ splits the diagonal scalar subgroup, compatibly with
all truncations. Consequently
\[
 L\simeq
 \left(\prod_{x\ne x_0}\GG_m\right)
 \times\prod_{x\in C(k)}U_x,
 \qquad U_x=1+t_x k[[t_x]].
\]
This equality is an equality of proalgebraic groups: each finite set of
points and each finite collection of jet orders occurs in a modulus.
The parameters $t_x$ are choices, irrelevant to the vanishing.
The big-Witt power-series isomorphism and its $p$-typical decomposition
\cite[Propositions 1.14 and 1.10, pp.~11 and 8]{HesselholtWitt} give
\[
 U_x\simeq\prod_{r\geq1,\ p\nmid r}W,
 \qquad W=\varprojlim_n W_n.
\]
Here $W$ is the infinite Witt group, not its universal cover.

Serre proves that $\GG_m$ has projective dimension one (\S7.5,
pp.~50--51), that $W$ has projective dimension one (\S8.6,
Corollary 3, p.~58), and that an abelian variety has projective dimension
at most one (\S9.2, Proposition 4, p.~61) \cite{Serre1960}.
For a product $T=\prod_i T_i$ of copies of $\GG_m$ or $W$, and a
\emph{quasi-algebraic} target $B$, Proposition 7 of \S3.4, p.~27, gives
\[
 \Ext^2_{\mathcal P}(T,B)
 =\varinjlim_{F\text{ finite}}
  \Ext^2_{\mathcal P}\left(\prod_{i\in F}T_i,B\right)=0.
\]
The proposition passes an inverse limit in the first variable to a
direct limit of Ext groups; its target hypothesis is essential and is
satisfied here. Finite products are finite direct sums. In particular,
the displayed vanishing holds for every elementary target, including
the finite groups $\ZZ/\ell$. Proposition 9 on the same page then
implies $\operatorname{pd}_{\mathcal P}T\leq1$ for \emph{all} targets.
Thus $\operatorname{pd}L\leq1$. The long exact Ext sequence for
$0\to L\to J_\infty\to J\to0$ now gives the result in every degree
at least two.
\end{proof}

\begin{proposition}[H2: the additive target; Proved]\label{h:additive}
For every $P\in\mathcal P$ and every connected quasi-algebraic $B$,
$\Ext^i_{\mathcal P}(P,B)=0$ for $i\geq2$.
In particular, $G=\GG_a$ is not the unresolved target suggested in
Part H's wording.
\end{proposition}
\begin{proof}
For an elementary source $\GG_a$, the assertion in degree two is
\cite[\S8.6, Corollary 5, p.~58]{Serre1960}. The other elementary
sources---finite groups, $\GG_m$, and abelian varieties---have
projective dimension at most one; for finite groups see \S4.4,
Corollary 2, pp.~34--35. Every quasi-algebraic group has a finite
composition series with these quotients (\S1.3, Corollary to
Proposition 7). The long exact sequence proves the assertion for a
quasi-algebraic source, and \S3.4, Proposition 7 passes to an arbitrary
proalgebraic source. Degrees at least three vanish by \S10.1,
Theorem 1. No inverse limit in the target has been used.
\end{proof}

\begin{proposition}[H3: augmentation cannot be omitted; Refuted comparison]
\label{h:augmentation}
The equality $\Ext^*(J_\infty,G)=H^*(K,G)$ is false in degree zero
for the connected $J_\infty$ above.
\end{proposition}
\begin{proof}
For $G=\ZZ/\ell$, a morphism $J_\infty\to G$ factors through a
connected $J_{\mathfrak n}$ and hence is zero. But
$H^0(K,\ZZ/\ell)=\ZZ/\ell$. On 43:80 the earlier object
$J'_{X/k}$ is explicitly an extension of $\ZZ$ by a connected group.
The boxed passage to $K$ must retain this augmentation; it cannot be
read as a statement about the connected object alone. The infinite
discrete group $\ZZ$ is not an object of Serre's $\mathcal P$.
Statement~\ref{h:vanishing} therefore settles the higher-Ext question
in the specified Serre category, without proving an equivalence of
Ext theories or the augmented all-degree comparison on another site.
\end{proof}

\subsection{The local-Jacobian lead and its topology}
\begin{proposition}[H4: an \'etale local-cohomology obstruction; Proved]
\label{h:local}
Let $A=k[[s,t]]$ and $\mathfrak m=(s,t)$. Then
\[
 H^2_{\mathfrak m,\mathrm{et}}(A,\ZZ/p)=0,
 \qquad
 H^3_{\mathfrak m,\mathrm{et}}(A,\ZZ/p)
 \simeq\bigoplus_{\substack{i,j>0\\p\nmid\gcd(i,j)}}(k,+)
\]
as $\mathbf F_p$-vector spaces. All other degrees vanish.
Thus the concentration used in SA19-H4 cannot be transferred from
Zariski cohomology to \'etale cohomology.
\end{proposition}
\begin{proof}
Quasi-coherent \'etale comparison and the \v Cech localization model
give $H^i_{\mathfrak m}(A,\GG_a)=0$ for $i\ne2$ and
\[
 E=H^2_{\mathfrak m}(A,\GG_a)
   =A_{st}/(A_s+A_t)
   =\bigoplus_{i,j>0}k e_{i,j},\qquad e_{i,j}=s^{-i}t^{-j}.
\]
All sums have finite support. Frobenius acts by
$F(ae_{i,j})=a^p e_{pi,pj}$. Partition the pairs into orbits
$p^r v$, $r\geq0$, where the two coordinates of $v$ are not both
divisible by $p$. On each orbit $F-1$ is injective: the last nonzero
coefficient of a proposed kernel element produces an uncancelled
coefficient one step later. Its cokernel is $(k,+)$, by the quotient
map
\[
 \sum_r a_r e_{p^r v}\longmapsto\sum_r a_r^{p^{-r}}.
\]
The map kills $F-1$, is surjective, and its kernel is exactly the
finite telescoping relations $F(ae_{p^r v})-ae_{p^r v}$.
The Artin--Schreier exact sequence now gives the asserted result.
The opened inputs are Stacks, Theorem 59.22.4, Section 59.63, and
Lemma 51.2.1 \cite{StacksEtaleQC,StacksAS,StacksLocalCech}.
\end{proof}

The ordinary functor $G\mapsto H^2_{\mathfrak m}(A,G)$ does not even
descend to $\mathcal P$ with its usual action on $\GG_a$:
Frobenius is an isomorphism in $\mathcal P$ and is not surjective on
$E$. If a projective representing object with this value and the
usual Artin--Schreier action were asserted, its Hom functor would make
$F-1:E\to E$ surjective, contradicting $e_{1,1}\notin(F-1)E$.
Even replacing $E$ by $\varinjlim_F E$ leaves a nonzero cokernel of
$F-1$: the orbits become indexed by $\ZZ$, still with finite support,
and the same telescoping map applies. This observation is conditional
on that particular attempted perfection, not a perfect-site comparison.

The reading explicitly describes 43:70--74 as Zariski theory.
Artin--Schreier is not exact on the Zariski site. The calculation
therefore refutes the proposed transfer and projectivity inference,
not an unspecified Zariski representability theorem. Completion from
$k[s,t]_{(s,t)}$ carries any alleged Artin--Schreier preimage to one
over $A$, so the obstruction also applies to that essentially finite
type example. The finite checks in \texttt{finite\_checks.py} retain
the extra target exponent and verify the obstruction for $p=2,3$;
the finite-support proof above establishes it for every $p$.

\begin{corollary}[H5: the perfect \'etale category; Proved]\label{h:perfect}
On the site $S=(\operatorname{Perf}/k)_{\mathrm{et}}$, for the
connected $J_\infty$ above and every quasi-algebraic target $G$,
\[
 \Ext^n_{\operatorname{Ab}(S)}(J_\infty,G)
   \simeq\Ext^n_{\mathcal P}(J_\infty,G)=0\qquad(n\geq2).
\]
Thus the higher vanishing in Part H holds in both stated categories,
with algebraic groups interpreted through perfection in those categories.
\end{corollary}
\begin{proof}
Suzuki's Proposition 3.8.1, p.~47 of the opened author version,
proves comparison for affine pro-quasi-algebraic sources and affine
quasi-algebraic targets. His Remark 3.8.4, pp.~55--56, extends it to
sources that are extensions of a perfected abelian variety by an
affine proalgebraic group, and arbitrary quasi-algebraic targets
\cite{SuzukiExt2022}. The sequence $0\to L\to J_\infty\to J\to0$
constructed above places $J_\infty$ in exactly that category.
Apply the comparison, then the proved projective-dimension bound.
\end{proof}

The comparison is not a formal consequence of a fully faithful exact
embedding: Suzuki explicitly states that proalgebraic realization on
$S$ is only left exact (p.~53). His proof passes through the perfect
pro-fppf site and the inverse-source limit theorem, Lemma 3.8.2,
pp.~48--53. No interchange of an arbitrary inverse limit in the target
is being asserted here.

For any actual exact sequence of sheaves
$0\to J_\infty\to J_{\mathrm{aug}}\to\ZZ\to0$, the same
higher vanishing follows. Indeed,
$\Ext^n_S(\ZZ,G)=H^n(S,G)=0$ for $n>0$: an \'etale cover of
the terminal object $\Spec k$ splits, so global sections are exact.
Apply the long exact sequence in the first variable. This does not
identify an unspecified augmented pro-object with ordinary
$H^*(K,G)$; the geometric argument and coefficient perfection still
have to be specified, especially for rational points in degree zero.

\begin{proposition}[H4b: a local Jacobian is not projective; Refuted claim]
Let $A_0=k[s,t]_{(s,t)}$ and let $x$ be its closed point.
The ordinary local Jacobian $J_x$ in
$\operatorname{Pro}(\operatorname{AbSch}_k)$, corepresenting
$G\mapsto H^2_{x,\mathrm{fppf}}(A_0,G)$ on commutative finite-type
$k$-group schemes, is not projective. In fact
\[
 \Ext^1_{\mathrm{naive}}(J_x,\ZZ/p)
 \simeq\bigoplus_{\substack{i,j>0\\p\nmid\gcd(i,j)}}(k,+).
\]
\end{proposition}
\begin{proof}
Existence and the stated universal property follow from To\"en's
construction and Definition 3.5, pp.~17--18, applied to the
codimension-two point of the smooth surface $\mathbb A^2_k$
\cite{Toen2026}. Let $f:J_x\to\GG_a$ correspond to
$e=s^{-1}t^{-1}\in E_0=H^2_x(A_0,\GG_a)$, and form the
pullback of the Artin--Schreier epimorphism along $f$:
\[
 Q=J_x\mathbin{\times}_{\GG_a,f,F-1}\GG_a,
 \qquad 0\longrightarrow\ZZ/p\longrightarrow Q
              \longrightarrow J_x\longrightarrow0.
\]
A splitting is exactly a morphism $g:J_x\to\GG_a$ with
$(F-1)g=f$, hence an element $h\in E_0$ with $(F-1)h=e$.
Completion carries this equation to the impossible equation in
$H^2_{(s,t)}(k[[s,t]])$ computed above. Thus the extension is
nonsplit, proving nonprojectivity in the stated category.

Local cohomology commutes with flat completion; since $E_0$ is
$\mathfrak m$-torsion, completion leaves it unchanged. Hence it is
the previously computed module $E$ with its Frobenius action.
To\"en's Lemma 3.9, pp.~21--22, gives
$\Ext^1_{\mathrm{naive}}(J_x,\GG_a)=0$.
The Artin--Schreier long exact sequence identifies the displayed
Ext group with $E/(F-1)E$, whose full orbit decomposition was proved
above. This works for every $p$, including $2$ and $3$.
\end{proof}

\paragraph{Scope of the modern geometric comparison.}
To\"en's Theorem 0.3 and Proposition 3.7 compare the full complex
of local Jacobians $J_*(X)$ with fppf cohomology for smooth $X$,
with $G$ affine if the cohomological degree is positive and unipotent
if it exceeds one. Section 3.2, p.~19, explicitly distinguishes
Yoneda Ext in pro-group schemes from Ext of all fppf sheaves;
the realization is not fully faithful in higher degrees.
Remark 3.10(2), p.~22, explains the missing restrictions in the old
formulation by its use of Zariski cohomology \cite{Toen2026}.
This is not a category-independent all-target comparison.

Cais--Otabe's Theorem 1.3, p.~3, and formula (1.8), p.~2,
prove
\[H^1_{\mathrm{fppf}}(U,G)=\varinjlim_{\mathfrak m}
\Ext^1(J_{\mathfrak m},G)\]
for every finite commutative group
scheme $G$, including nonreduced ones \cite{CaisOtabe2026}.
Their p.~3 discussion notes that Deligne's 1974 letter asserted
arbitrary commutative $G$ but left details in a two-variable
effaceability argument. Their finite-target theorem is therefore
not used as an unqualified all-degree comparison.

\begin{proposition}[H3b: augmentation alone does not repair ordinary degree zero; Refuted claim]
On the perfect \'etale site, no source object $P$ can give
isomorphisms
\[
 \Hom(P,G^{\mathrm{perf}})\simeq H^0(k(t),G)
\]
natural in all commutative algebraic $k$-groups $G$, where the
right side means ordinary rational points. This applies also to
any proposed augmented $J_{\mathrm{aug}}$.
\end{proposition}
\begin{proof}
Take $G=\GG_a$ and its homomorphism $F:x\mapsto x^p$.
The induced endomorphism of $G^{\mathrm{perf}}$ is an isomorphism,
so it acts bijectively on $\Hom(P,G^{\mathrm{perf}})$ for every
$P$. On $H^0(k(t),\GG_a)=k(t)$ it is not surjective: if
$h^p=t$, the valuation at $t=0$ gives $p\,v_0(h)=1$, a
contradiction. Naturality would conjugate these two maps, which
is impossible. Thus both the augmentation and the perfection of
the geometric/coefficient argument must be specified before an
all-degree comparison can be asserted. This does not affect
higher Ext vanishing in the precisely stated categories.
\end{proof}
